\label{mdg:chap:barbero}

La transición de una hexada a una octada combina tres estructuras de tipos
distintos: una incidencia finita, dos series de Lambert asociadas a los
canales \(q_\pm\) y una identificación con el parámetro real de la acción de
Holst. Las tres se articulan mediante la incidencia, el funcional angular y
el diccionario areal hacia la geometría de la conexión de
Ashtekar--Barbero desarrollado en esta misma residencia.

\section{Cardinalidades de incidencia 90 y 120}

Sea \(H\) un conjunto de seis puntos y sea \(O\) un conjunto de ocho puntos
que contiene una copia distinguida de \(H\). Consideremos
\[
 \mathcal F_6=H\times\binom H2,
 \qquad
 \mathcal F_8=O\times\binom H2.
\]
La inclusión distinguida \(H\hookrightarrow O\) induce
\[
 \iota_{6,8}:\mathcal F_6\hookrightarrow\mathcal F_8,
 \qquad
 (h,\{h_1,h_2\})\longmapsto(h,\{h_1,h_2\}).
\]
Este morfismo de incidencias es el antecedente combinatorio de los
cardinales \(90/120\); no es el extractor dodecafásico ni una serie angular.

\begin{theorem}[Cardinalidades de la inclusión]
\[
 |\mathcal F_6|=6\binom62=90,
 \qquad
 |\mathcal F_8|=8\binom62=120.
\]
La diferencia normalizada por veinticuatro coordenadas es
\[
 \frac{90-120}{24\binom62}=-\frac1{12}.
\]
\end{theorem}
\begin{proof}
Cada elemento de la primera coordenada puede combinarse con cualquiera de los
quince pares de \(H\). El principio multiplicativo da las dos cardinalidades;
la última identidad resulta de sustituir \(\binom62=15\).
\end{proof}

\begin{figure}[htbp]
\centering
\begin{tikzpicture}[
  font=\small,
  setnode/.style={circle,draw,minimum size=6mm,inner sep=0pt,
    font=\scriptsize\bfseries},
  box/.style={draw,rounded corners=1.5mm,align=center,
    text width=5.0cm,inner xsep=6pt,inner ysep=5pt},
  arr/.style={-{Stealth[length=2.2mm]},line width=.7pt}]

  % Inclusión geométrica. Las flechas parten del contorno y no atraviesan nodos.
  \node[setnode,draw=hmtblue,fill=hmtlightblue] (h1) at (-2.3,1.25) {1};
  \node[setnode,draw=hmtblue,fill=hmtlightblue] (h2) at (-1.1,1.9) {2};
  \node[setnode,draw=hmtblue,fill=hmtlightblue] (h3) at (.1,1.25) {3};
  \node[setnode,draw=hmtblue,fill=hmtlightblue] (h4) at (.1,-.15) {4};
  \node[setnode,draw=hmtblue,fill=hmtlightblue] (h5) at (-1.1,-.8) {5};
  \node[setnode,draw=hmtblue,fill=hmtlightblue] (h6) at (-2.3,-.15) {6};
  \node[setnode,draw=hmtgreen,fill=hmtlightgreen] (o7) at (-3.2,.55) {7};
  \node[setnode,draw=hmtgreen,fill=hmtlightgreen] (o8) at (1.0,.55) {8};
  \node[draw=hmtblue,dashed,rounded corners=4mm,
    fit=(h1)(h2)(h3)(h4)(h5)(h6),inner sep=4mm,
    label={[hmtblue,font=\bfseries]above:{hexada \(H\), \(|H|=6\)}}] (H) {};
  \node[draw=hmtgreen,line width=.8pt,rounded corners=6mm,
    fit=(H)(o7)(o8),inner sep=5mm,
    label={[hmtgreen,font=\bfseries]below:{octada \(O\), \(|O|=8\), \(H\subset O\)}}] (O) {};

  \node[box,draw=hmtblue,fill=hmtlightblue] (f6) at (5.3,1.35)
    {Incidencias definidas sobre \(H\)\\[1mm]
     \(\mathcal F_6=H\times\binom H2\),\quad
     \(|\mathcal F_6|=6\binom62=90\)};
  \node[box,draw=hmtgreen,fill=hmtlightgreen] (f8) at (5.3,-1.0)
    {Extensión de incidencias a \(O\)\\[1mm]
     \(\mathcal F_8=O\times\binom H2\),\quad
     \(|\mathcal F_8|=8\binom62=120\)};
  \draw[arr,hmtblue] (O.east) -- ++(.45,0) |- (f6.west);
  \draw[arr,hmtgreen] (O.east) -- ++(.45,0) |- (f8.west);
  \draw[arr,hmtgold] (f6.south) -- node[right,align=left,font=\footnotesize]
    {inclusión \(\iota_{6,8}\)} (f8.north);

  \node[box,draw=hmtgold,fill=hmtlightgold,text width=10.9cm]
    (def) at (1.45,-3.80)
    {Defecto combinatorio orientado\\[-.2mm]
     con la normalización declarada:\\[1mm]
     \(\displaystyle
       \frac{|\mathcal F_6|-|\mathcal F_8|}{24\binom62}
       =\frac{90-120}{24\cdot15}=-\frac1{12}.\)};
  \draw[arr,hmtgold] (f8.south) -- (f8.south |- def.north);

  \node[align=center,text=hmtgray,font=\footnotesize,text width=11.2cm]
    at (1.45,-5.50)
    {Esta identidad pertenece al mapa de incidencias. No identifica el
     extractor dodecafásico, el semigrupo de torres ni el funcional angular.};
\end{tikzpicture}

\caption{Inclusión de un subconjunto distinguido de seis elementos en uno de
ocho elementos y conteos de incidencia. Los cardinales \(90\) y \(120\) no fusionan el extractor
dodecafásico, el funcional angular ni el semigrupo de torres.}
\label{mdg:fig:incidencia-seis-ocho-rev9}
\end{figure}

Los índices iniciales de la jerarquía angular quedan así asociados a dos
espacios de incidencia: 90 configuraciones sobre la hexada y 120 sobre la
octada ambiente. El cociente \(-1/12\) es una identidad combinatoria de esta
normalización.

\section{Series de Lambert asociadas a los 2 sectores de incidencia}

Para \(0<q<1\), definimos simultáneamente
\[
 S_{90}(q)=\frac{q^{90}}{1-q^{270}},
 \qquad
 S_{120}(q)=\frac{q^{120}}{1-q^{360}},
 \qquad
 \Delta S_m=S_m(q_-)-S_m(q_+).
\]
Las expansiones geométricas
\[
 \Delta S_{90}=\sum_{j\geq0}s_{90+270j},
 \qquad
 \Delta S_{120}=\sum_{j\geq0}s_{120+360j}
\]
sitúan ambos términos dentro del semigrupo global del
\cref{mdg:chap:torres}.

El funcional angular asociado a la transición \(6\to8\) es
\[
 K_{6\to8}=12\Delta S_{90}-\Delta S_{120},
\]
y su precedencia ha quedado fijada por la cadena fundamental dodecafásica:
los doce sectores y la única costura orientada producen primero el peso
\(12:-1\). El coeficiente de \((1-q)^{-1}\), igual a \(1/24\), se calcula después como salida exacta;
la ecuación diofántica que sigue certifica el par ya generado y no lo
selecciona.
La coordenada angular resultante es
\[
\Gamma_{6\to8}
=\frac{\pi A\Cstar}{180}+K_{6\to8}.
\]
En el punto HMT,
\[
\begin{aligned}
 \Delta S_{90}&=2.95169439297457621139847987861\times10^{-4},\\
 \Delta S_{120}&=1.96835112502951720093738706217\times10^{-5},\\
 \frac{\pi A\Cstar}{180}&=0.270485322934140078465586458790\ldots,
\end{aligned}
\]
y, por tanto,
\[
 \boxed{\Gamma_{6\to8}
 =0.27400767269445927474725526077398041940\ldots.}
\]
Las tablas calculadas con la coordenada angular conjugada \(C^*\) redondeada a quince decimales dan
\(0.2740076726944593\) en aritmética de doble precisión; el valor anterior
es la evaluación multiprecisión de la fase regional completa.

\section{Certificado diofántico posterior de los pesos}

La identidad
\[
 \lim_{q\to1^-}(1-q)\frac{q^m}{1-q^n}=\frac1n
\]
asocia a la combinación \(aS_{90}-bS_{120}\), antes de su evaluación
en los canales conjugados, el coeficiente de \((1-q)^{-1}\)
\[
 \frac a{270}-\frac b{360}.
\]

\begin{theorem}[Pesos enteros con coeficiente de polo \texorpdfstring{\(1/24\)}{1/24}]
Las soluciones enteras positivas de
\[
 \frac a{270}-\frac b{360}=\frac1{24}
\]
son
\[
 (a,b)=(12+3t,1+4t),\qquad t\in\ZZ_{\geq0}.
\]
El par \((12,1)\) es la única solución mínima coordenada a coordenada y la
única de norma \(\ell^1\) mínima.
\end{theorem}
\begin{proof}
Multiplicar por \(1080\) produce \(4a-3b=45\). Módulo tres se obtiene
\(a\equiv0\pmod3\); escribiendo \(a=3u\), resulta \(4u-b=15\). La
positividad implica \(u=4+t\), \(b=1+4t\), con \(t\geq0\). Por tanto,
\(a+b=13+7t\), lo que demuestra la minimalidad.
\end{proof}

La primitividad \(\gcd(a,b)=1\) no selecciona por sí sola el par mínimo, pues
existen otras soluciones primitivas. La vuelta fundamental del TPK ya ha
producido el peso \(12:-1\); el coeficiente de polo \(1/24\), la positividad y la
minimalidad forman su certificado aritmético posterior. Ninguno de estos
objetos se deduce de la evaluación decimal de \(\Gamma_{6\to8}\).

\newpage
\section{Monotonía del funcional respecto de la coordenada angular conjugada}

Para \(m\in\{90,120\}\),
\[
 S_m'(q)=\frac{m q^{m-1}(1+2q^{3m})}{(1-q^{3m})^2}>0.
\]
Además,
\[
 \frac{\dd q_-}{\dd C}=\frac\pi{180}q_-,
 \qquad
 \frac{\dd q_+}{\dd C}=-\frac\pi{180}q_+.
\]

\begin{proposition}[Monotonía estricta]
Para \(1^\circ\leq C\leq3^\circ\), la función
\(C\mapsto\Gamma_{6\to8}(A,C)\) es estrictamente creciente. En consecuencia,
todo valor de su imagen tiene a lo sumo una preimagen en ese intervalo.
\end{proposition}
\begin{proof}
La derivada de cada diferencia \(\Delta S_m\) es positiva. En el intervalo
indicado se obtiene
\(\Delta S'_{120}<5.35\times10^{-4}\), mientras que
\(\pi A/180>0.1273\). Por tanto,
\(\Gamma'(C)>0.1267\). En \(C=\Cstar\),
\[
 \Gamma'(\Cstar)=0.1328995105618907\ldots.
\]
\end{proof}

La inyectividad local excluye una segunda coordenada angular próxima con la
misma evaluación del funcional. Este control no construye ni selecciona
\(C^*\): la coordenada angular conjugada ya ha sido obtenida en la rama regional anterior.
La proposición sólo cuantifica la sensibilidad de la evaluación descendente
una vez fijados \(A\) y \(C^*\); no autoriza a invertir
\(\Gamma_{6\to8}\) ni un valor físico de \(\gamma_{\mathrm{BI}}\) para
redefinirlos.
